\documentclass[10pt,a4paper]{article}
\usepackage[T1]{fontenc}\usepackage[utf8]{inputenc}
\usepackage{lmodern,amsmath,amssymb,booktabs,array,longtable,graphicx}
\usepackage[margin=23mm]{geometry}\usepackage{microtype,xcolor,hyperref}
\hypersetup{colorlinks=true,linkcolor=blue!45!black,urlcolor=blue!45!black,pdftitle={WestQuant: System Evaluation 1--4},pdfauthor={David Vesterlund}}
\setlength{\emergencystretch}{3em}\setlength{\parskip}{3pt}
\newcommand{\code}[1]{\texttt{#1}}
\title{System Evaluation 1--4:\\Methods, Executed Results and Evidence Limits}
\author{David Vesterlund\\\small WestQuant, Vesterlund Ventures, Sweden}
\date{10 October 2026}
\begin{document}\maketitle
\noindent This dossier records the four system-evaluation pathways. Steps 1, 3 and 4 are executed at their declared software scope. Step 2 has an executed synthetic molecular-spin control, with experimental physical validation still pending. The content is also integrated into the main manuscript and Supplement. No improved material lifetime, hardware result or prediction confidence interval is claimed.
\begin{table}[htbp]\centering\small\begin{tabular}{p{.62\linewidth}r}\toprule Quantity & Result\\\midrule
Profiles / channel-age points & 13 / 73\\
Processed fidelity values & 90\\
Maximum full Choi discrepancy & $1.310\times10^{-14}$\\
Maximum six-state output discrepancy & $1.183\times10^{-14}$\\
Maximum conditional average-fidelity discrepancy & $3.997\times10^{-15}$\\
Maximum acceptance / processed-value discrepancy & $0$ / $0$\\
Declared numerical gate & $1.000\times10^{-8}$\\\bottomrule\end{tabular}\caption{Step 1: matched numerical and processed-data comparison. Twelve profiles are supplied software controls; one profile is an experiment-constrained point.}\label{tab:systemcomparison}\end{table}

\section{Additional system evaluation: complete executed protocol}
\label{sec:systemprotocol}
Steps 1, 3 and 4 have executed results at their stated software scope. Step 2 has an executed synthetic second-family case, whereas its experimental physical objective remains blocked. The four selected records share evaluation-source fingerprint
\par\noindent\begin{minipage}{\linewidth}\centering\small\texttt{d8fc6d2ecb580e0ba1c8881ffbd542d087c}\\\texttt{493b7621cb89509948e595211ec92}.\end{minipage}\par
Their common lab-source fingerprint is
\par\noindent\begin{minipage}{\linewidth}\centering\small\texttt{9c4490acf2e5a42e8e10007989eeb94b5e4}\\\texttt{06b17f5d187ca310847581ee2edb7}.\end{minipage}\par
These fingerprints identify source content, not hardware or an external scientific endorsement. Steps 2 and 3 were confirmed with the published implementation after their initial runs; the initial outputs remain archived. Table generation asserts the shared fingerprint and expected numerical/status gates without overwriting the original run records. The configuration digest and UTC timestamp of each selected record are in \code{evidence/system\_evaluation\_claim\_map.json}. Generated table files are under \code{evaluation/paper/}.

\subsection{Step 1: independent reference equations and matched inputs}
The reference \code{evaluation/wqeval/reference.py} imports NumPy/SciPy without importing \code{memorylab}. For column vectorization and Hamiltonian $G=\mathrm{diag}(0,\delta)$ in Hz, it uses
\[
\mathcal L=-2\pi i(I\otimes G-G^{\mathsf T}\otimes I)
+\sum_J\left[J^*\otimes J-\tfrac12 I\otimes J^\dagger J
-\tfrac12(J^\dagger J)^{\mathsf T}\otimes I\right].
\]
The jump operators are $\sqrt{\Gamma_1(1-p_e)}\,|0\rangle\langle1|$, $\sqrt{\Gamma_1p_e}\,|1\rangle\langle0|$, and $\sqrt{\Gamma_\phi/2}\,Z$. The stage superoperator is $\exp(t\mathcal L)$. Write and read stages include their specified depolarizing maps, and storage adds the declared network-activity dephasing. The reference constructs $C=\sum_{a,b}|a\rangle\langle b|\otimes\mathcal E(|a\rangle\langle b|)$ by matrix-unit propagation, including an orthogonal erasure outcome for the full channel. It computes six Pauli-eigenstate outputs and their mean target fidelity independently of the lab's Kraus construction.

The four joint parameter choices $(\Gamma_1,\delta,p_e)$ are $(0,0,0)$, $(0.05,123,0.1)$, $(10,-500,0.5)$ and $(100,1000,0.9)$, with rates in s$^{-1}$ and detunings in Hz. Each is combined with storage temperatures 4.2, 10 and 80 K and ages $[0,10^{-6},10^{-4},10^{-3},10^{-2},10^{-1}]$ s. Write/read stages remain at 4.2 K, last $10^{-4}/2\times10^{-4}$ s and use acceptance 0.8/0.7 and depolarization 0.01/0.02. The pure-dephasing rate is $0.01+0.001T+0.001(T/10)^3$ s$^{-1}$, erasure rate 0.01 s$^{-1}$, and storage activity contributes $10\times0.002=0.02$ s$^{-1}$. Rate validity is declared over 1--300 K. These are chosen software profiles, not fitted material parameters. The measured Teller point supplies the thirteenth profile and its sole supported age of 18 $\mu$s at 3 K, giving $12\times6+1=73$ channel-age points.

The two experimental input files are \code{4.csv} and \code{8a.csv} from record 14929997:v1. Their SHA-256 hashes, original download URLs and data-license metadata are pinned in \code{evaluation/data/teller\_sources.json}. Both parsers use the same checked bytes. The path file yields 30 values and mean 0.9510824228990328; the time-bin file yields 60 values and mean 0.9140479475068066. The discrepancy is zero for every processed value. The comparison threshold is $10^{-8}$ for the full Choi matrix, six states, average fidelity, acceptance and processed values. This threshold is numerical, not an experimental accuracy margin.

Both timing arms are warmed up before five repeats, with their order randomized by seed 20261010. All 73 channel constructions are included in each repeat. \code{construction\_timing.csv} retains every measured duration and actual order. The arms do not perform equal amounts of metadata, export, validation or metric-construction work. File IO, source retrieval, network events, fits and electronic structure are excluded. Native thread environment variables are recorded but were not constrained for step 1. The medians cannot be interpreted as a fair end-to-end speedup or engineering-effort reduction. \code{related\_tools\_audit.json} distinguishes executed comparison from documentary context; AiiDA, PySCF, PyCCE, QuTiP and SeQUeNCe were not separately timed here. Existing AiiDA provenance and caching are not presented as new contributions.

\subsection{Step 2: synthetic control, curve gate and physical stop}
The model/curve comparison requires the observable \code{complex\_pure\_dephasing\_coherence}, explicit exclusion of population relaxation, source hashes, curve reviewer, normalization and observable mapping. A measured case additionally requires reviewed model-parameter provenance, a nonhypothetical model and exact equality of full declared experimental conditions between curve and model. The fixed-model runner accepts \code{exact\_spin\_bath} or optional \code{pycce}; it fits zero parameters. The initial curve value must be unity, magnitudes must lie in the unit disk and a strictly increasing finite grid must start at zero. No silent clipping of noisy or unnormalized measurements is provided.

The executed control has one $^{13}$C nuclear spin $I=1/2$, central logical projections $\pm1/2$, laboratory $A_{zz}=4000$ Hz and $\nu_z=214163.9829$ Hz, zero other couplings, 4.2 K and zero pulses. Forty-one equally spaced ages span 0--1 ms. An independent closed-form Gibbs-initialized FID produces the curve; the finite Hamiltonian solver supplies the comparison. Every complex residual is divided by the chosen scale 0.01, and the maximum must be at most three. The observed maximum is $3.7192473\times10^{-12}$. This chosen scale is not an experimental SEM, and the gate has no calibrated statistical false-positive rate. The 75--100 $\mu$s crossing bracket differs from the earlier 101-point control's 90--100 $\mu$s bracket solely because the grids differ; both are sampled numerical brackets and both detect revivals.

All write/store/read temperatures are hypothetically 4.2 K. The supplied interface has acceptance 0.8/0.7, depolarization 0.01/0.02 and durations 100/200 $\mu$s, with the independent rate models retained from the example. Export requires \code{--allow-hypothetical}. At 500 $\mu$s, the provider-neutral two-node workload has 10,000 trials, herald probability 0.1, 1 ms interval, Bell-fidelity threshold 0.8 and seed 20261010. It yields 1020 heralds and 313 retrieved/quality-qualified pairs. Conditional Bell fidelity is 0.9559176317211326; the retrieved-fraction Wilson interval is [0.0280630600,0.0348969001]. The simulated qualified rate is 31.3 s$^{-1}$ under these supplied assumptions, not measured network performance. This workload is not the earlier 2000-trial SeQUeNCe control.

The separate pending descriptor targets a molecular reference at DOI 10.1039/C8SC04122A. Running it returns \code{blocked\_missing\_reviewed\_data} and exit code two, with no physical validation and no network export. Required inputs are a reviewed normalized curve with measurement uncertainties, fixed Hamiltonian and tensor frames, field/temperature/isotopes/pulses, separation of relaxation from pure dephasing, and independent optical-interface review. Reported molecular phase-memory time alone does not satisfy these requirements. Thus no second-family experimental reproduction or new-material lifetime is supplied by step 2.

\subsection{Step 3: individual fault interventions}
The suite creates local copies of inputs and explicitly unsafe bypasses without modifying production guards. Each invalid case has a valid positive control that must succeed. Rejection alone would be insufficient if every valid input also failed. All seven positive controls and all seven guarded rejections pass; the deliberate bypasses accept all seven cases.
{\small\begin{longtable}{p{.37\linewidth}p{.16\linewidth}p{.38\linewidth}}\toprule Constructed fault & Guarded result & Accepted unsafe-bypass consequence\\\midrule\endhead
Operating-point age extrapolation & Blocked & $10\times$ unsupported age\\
Operating-point temperature extrapolation & Blocked & 1 K unsupported shift\\
Total coherence rate reused as pure dephasing & Blocked & 16.667\% conditional $T_2$ error\\
Unreviewed noise independence & Blocked & One unsupported declaration\\
Hypothetical channel without opt-in & Blocked & Two hypothetical age samples\\
Absolute-age map applied incrementally & Blocked & 0.097247 maximum element error\\
Modified source bytes & Blocked & $1000\times$ forged source value\\\bottomrule\caption{Step 3: all seven valid positive controls pass; all invalid cases are rejected. Consequences have distinct units and are not aggregated into an accuracy score.}\label{tab:systemfaults}\end{longtable}}


The age bypass changes the declared sole point age before evaluating a tenfold longer age, thereby illustrating an unsupported extrapolation rather than measuring its physical error. The temperature bypass changes the declared point temperature by 1 K, again with no experimentally known extrapolation error. Noise independence is counted as an unsupported declaration; evidence promotion counts two age samples. The rate-definition example has $\Gamma_1=10$ s$^{-1}$ and correct independent $\Gamma_\phi=20$ s$^{-1}$, giving 40 ms. Treating total coherence rate 25 s$^{-1}$ as pure dephasing gives 33.333 ms instead. The modified-source case changes a chosen value from 0.001 to 1 s; it exercises the new evaluation source contract, whereas the other six cases exercise existing lab guards.

In the incremental-application control, the 1 ms complete channel includes write/read depolarization as well as storage; applying it twice produces maximum density-element discrepancy 0.0972467329 relative to applying it once. A separate units control converts 1 ms correctly to 0.001 s and contrasts it with an omitted conversion to 1 s; its maximum element discrepancy is 0.1322008762. These consequences have different units and scientific meanings. Neither their average nor seven out of seven represents a general material-prediction accuracy or real-world researcher error rate. Equivalent guards may be implemented in other frameworks.

\subsection{Step 4: workload, timing scope and robustness}
The native runner workload has Markov and one-spin exact-bath models, detunings $[0,123,500]$ Hz, two replicates and six ages $[0,10^{-4},5\times10^{-4},10^{-3},2\times10^{-3},5\times10^{-3}]$ s. Replicates have distinct deterministic seed streams; they are not independent hardware measurements or stochastic bath realizations in this deterministic workload. The Cartesian count is $2\times3\times2=12$ jobs per study. Master seed is 20261010. Three repeats independently randomize worker order among 1, 2 and 4. Worker count is excluded from physics identity, allowing a direct consistency check across scheduling choices.

Each cold study runs in a separate output directory and an isolated measurement process. OMP, MKL, OpenBLAS and NumExpr thread limits are one in those benchmark subprocesses. Per-job timeout is 60 s, driver timeout 600 s and max-run budget 24. The measured interval includes process startup, actual channel construction and artifact IO. The warm interval reuses terminal-result checkpoints; it is not electronic-structure memoization. Nine cold studies complete 108 jobs, and nine warm passes reuse 108 results. The scientific digest, excluding durations and output paths but including configuration identities and channel data, is identical across all nine studies.
\begin{table}[htbp]\centering\small\begin{tabular}{r r r r}\toprule Workers & Cold median [min,max] (s) & Resume median (s) & Median jobs/s\\\midrule
1 & 3.8674 [3.8089, 3.9101] & 0.01814 & 3.103\\
2 & 1.9781 [1.9165, 2.0665] & 0.01741 & 6.066\\
4 & 1.1782 [1.1299, 1.2286] & 0.01817 & 10.185\\\bottomrule\end{tabular}\caption{Step 4: three local cold studies and three checkpoint-reuse passes per worker count, with 12 jobs per study. Bounds are observed minima/maxima, not confidence intervals.}\label{tab:systemperformance}\end{table}


The reported bounds are observed ranges from three repeats, not confidence intervals. Separate resource measurements use process usage counters: parent maximum RSS is 34.578125--37.078125 MiB and the largest terminated child's maximum is 63.238281--66.414063 MiB. The metric is not peak total memory across concurrent workers. Child CPU seconds and each measurement's raw counters remain in the retained summary. Performance applies to this Linux environment and small supplied models, not general CCE, electronic structure, all elements or cluster-scale workloads.

Three additional robustness probes are distinct from the 108 cold jobs. A profile without write temperature produces \code{missing\_data}. A real worker process with a $10^{-6}$ s timeout produces \code{timeout}; this deadline intentionally tests termination and is not a solve-time estimate. A precompleted valid job followed by the full study completes all 12 jobs with exactly one checkpoint reused. Thus robustness adds one missing-data worker, one timeout worker and 12 successful calculations across the partial/restarted study. No forced-kill during an atomic write, power-loss, corruption or concurrent cache-writer experiment was performed. Failed terminal checkpoints are preserved as failures, not retried or counted as successful physics results.

\subsection{Core fresh-checkout CI and traceable work units}
GitHub Actions run 38051481969 at commit \code{de52b32f0c34a05cf4f6a0c01bca6ee087e6cb67} uses Ubuntu 24.04 and Python 3.12.15. It installs the core editable package with NumPy 2.5.3 and SciPy 1.18.1, passes seven evaluation tests, automatically retrieves/checks both pinned experimental files, and reruns the matched comparison. Its maximum full-Choi discrepancy is $1.3214015\times10^{-14}$ versus the $10^{-8}$ gate. \code{evidence/core\_ci\_job.log} and \code{core\_ci\_verification.json} preserve the log and summary; the log hash is checked. This core CI result is additional environment evidence, not full optional-backend reproduction or colleague-led independent scientific review. CSV bytes are preserved and automatic Git line-ending conversion is disabled to protect source hashes.

The original 75 tests plus seven additional evaluation tests constitute 82 distinct local passing tests. Repeating controls, CI tests, numerical ages or Monte Carlo trials does not increase that count. Work units remain separate: 73 channel-age comparisons, 90 processed data values, 41 synthetic coherence ages, 10,000 provider-neutral network trials, seven selected invalid contracts with seven positive controls, 108 cold worker jobs and 108 checkpoint reuses. No new hardware experiment or ML training run is counted. Steps 5 and 6 remain the author's and colleagues' full independent reproduction and scientific/release review tasks.

\section{Reproduction and data locations}
From the repository root, install the core package with \code{python -m pip install -e ./lab}. Commands and full source contracts are documented in \path{evaluation/README.md}. First-run step 1 retrieves its two pinned Zenodo inputs and verifies SHA-256 before parsing. Original third-party caches with unresolved redistribution permissions are absent from the public clone. Optional physics backends and full original article reproduction use the separate setup in \path{REPRODUCIBILITY.md}. Generate these paper tables and the evidence claim map with \code{python} \path{scripts/document_system_evaluation.py}.

{\raggedright The four selected run summaries are \path{evaluation/results/step1-executed/summary.json}, \path{evaluation/results/step2-confirmation/summary.json}, \path{evaluation/results/step3-confirmation/summary.json} and \path{evaluation/results/step4-executed/summary.json}. The missing-data physical case is \path{evaluation/results/step2-physical-pending/summary.json}. The machine-readable evidence index is \path{evidence/system_evaluation_claim_map.json}. Full per-case tables retain units and every timing repeat; the CI log and verification record are under \path{evidence/}.\par}

Source and article materials: \url{https://github.com/VesterlundCoder/memory-discover-lab}. Colleague-led independent reproduction and scientific/release review (steps 5 and 6) remain pending.
\end{document}
